\documentclass[11pt,a4paper]{article}
\usepackage[margin=2.2cm,landscape]{geometry}
\usepackage[T1]{fontenc}
\usepackage{lmodern}
\usepackage{booktabs}
\usepackage{longtable}
\usepackage{array}
\usepackage{amssymb}
\usepackage[hidelinks]{hyperref}

% variable names may contain underscores; \var prints them verbatim in typewriter and may break at underscores
\usepackage{url}
\makeatletter\g@addto@macro\UrlBreaks{\do\_}\makeatother
\DeclareUrlCommand\var{\urlstyle{tt}}
\newcolumntype{L}[1]{>{\raggedright\arraybackslash}p{#1}}
\newcommand{\yes}{\checkmark}
\newcommand{\no}{--}

\title{CGPS variables potentially used to define the outcome}
\author{}
\date{Compiled from \texttt{header\_BL.xlsx}, \texttt{header\_FU.xlsx} and the aggregate report of \texttt{cgps\_stats.py}}

\begin{document}
\maketitle

\noindent This document lists every variable in the CGPS tables that could be used to define the outcome
of the risk prediction model, or eligibility for it. The CT volumes are not listed. BL and FU denote the
baseline and follow-up tables; a tick means that the variable exists in that table. The role column refers
to the chosen design (Section~\ref{sec:design}): \emph{outcome}, \emph{eligibility}, \emph{check}
(sensitivity or noise check), \emph{covariate}, or \emph{not used}. Counts are aggregates from the report
on 2\,116 patients with both scans.

\section{Segments}
\label{sec:segments}

Segment-level variables are named \var{ct_cag_<segment>_<item>}. The 17 segments are listed below.
Only the trunk segments of the LAD, CX and RCA enter the chosen design; the LM enters as an eligibility
criterion for the LAD and CX.

\begin{longtable}{L{2.2cm} L{4.6cm} L{2.0cm} L{2.4cm} L{9.5cm}}
\toprule
Code & Segment & Vessel & Design role & Note \\
\midrule
\endhead
\var{lm}   & Left main                     & LM  & eligibility & Must be plaque-free for LAD and CX units \\
\var{plad} & Proximal LAD                  & LAD & outcome     & Trunk; 446 of 1\,003 plaque-free at BL \\
\var{mlad} & Mid LAD                       & LAD & outcome     & Trunk \\
\var{dlad} & Distal LAD                    & LAD & outcome     & Trunk \\
\var{d1}   & First diagonal                & LAD & not used    & Side branch \\
\var{d2}   & Second diagonal               & LAD & not used    & Side branch \\
\var{pcx}  & Proximal circumflex           & CX  & outcome     & Trunk \\
\var{mcx}  & Mid circumflex                & CX  & outcome     & Trunk \\
\var{dcx}  & Distal circumflex             & CX  & outcome     & Trunk \\
\var{om1}  & First obtuse marginal         & CX  & not used    & Side branch \\
\var{om2}  & Second obtuse marginal        & CX  & not used    & Side branch \\
\var{prca} & Proximal RCA                  & RCA & outcome     & Trunk \\
\var{mrca} & Mid RCA                       & RCA & outcome     & Trunk; most often unevaluable (motion) \\
\var{drca} & Distal RCA                    & RCA & outcome     & Trunk; ends at the PDA origin \\
\var{pda}  & Posterior descending artery   & RCA/CX & not used & Side branch; origin depends on dominance \\
\var{pla}  & Posterolateral branch         & RCA/CX & not used & Side branch; origin depends on dominance \\
\var{im}   & Ramus intermedius             & --  & not used    & Present in a minority; often coded unevaluable when absent \\
\bottomrule
\end{longtable}

\subsection{Per-segment items}

\begin{longtable}{L{4.3cm} c c L{6.2cm} L{2.2cm} L{7.0cm}}
\toprule
Variable & BL & FU & Values & Role & Note \\
\midrule
\endhead
\var{ct_cag_<s>_stenosis} & \yes & \yes &
  Text: none, minimal, mild, moderate, severe, occlusion, \var{uneval_*} &
  outcome, eligibility &
  Decoded to grades 0--5; \var{uneval_*} treated as missing. Empty at BL for 1\,094 patients (not read). \\
\var{ct_cag_<s>_plaque_type} & \yes & \yes &
  Text: calcified, mostly\_calcified, mostly\_noncalcified, noncalcified; empty when no plaque &
  check &
  Composition of the new plaque; plaque presence is taken from the stenosis grade. \\
\var{ct_cag_<s>_feat_spotcalc} & \yes & \yes & 0/1 & not used &
  High-risk feature; recorded only with plaque; mostly noise (opposite changes exceed new findings). \\
\var{ct_cag_<s>_feat_napring} & \yes & \yes & 0/1 & not used & High-risk feature; rare. \\
\var{ct_cag_<s>_feat_lap} & \no & \yes & 0/1 & not used & Low-attenuation plaque; 0 for every patient at FU. \\
\var{ct_cag_<s>_feat_posrem} & \no & \yes & 0/1 & not used & Positive remodelling; FU only, so no change can be computed. \\
\var{ct_cag_<s>_eval} & \yes & \no & Numeric code (0, 1, 5, 12--14) & check &
  Evaluability; meaning of the codes unknown, to be confirmed with the data manager. \\
\var{ct_cag_<s>_medis_egnet} & \yes & \no & Numeric & not used & Suitability for Medis plaque analysis. \\
\var{ct_cag_<s>_medis_finished} & \yes & \no & Numeric & check &
  Medis analysis done ($\approx$200 patients); exports may contain plaque positions along the centreline. \\
\bottomrule
\end{longtable}

\subsection{Stenosis grading}

The grade labels match the SCCT categories \cite{leipsic2014,cury2016}. The percentage ranges below are those
of the SCCT guideline and are assumed, not confirmed, for the CGPS reads.

\begin{center}
\begin{tabular}{l c l l}
\toprule
Label & Grade & Diameter stenosis (SCCT) & Category \\
\midrule
none      & 0 & 0\,\%        & no plaque \\
minimal   & 1 & 1--24\,\%    & non-obstructive \\
mild      & 2 & 25--49\,\%   & non-obstructive \\
moderate  & 3 & 50--69\,\%   & obstructive \\
severe    & 4 & 70--99\,\%   & obstructive \\
occlusion & 5 & 100\,\%      & obstructive \\
\bottomrule
\end{tabular}
\end{center}

\section{Vessel-level summaries}

Computed by the CGPS readers from the segment reads, for \var{<v>} $\in$ \{\var{lad}, \var{cx}, \var{rca}\}.

\begin{longtable}{L{4.3cm} c c L{6.2cm} L{2.2cm} L{7.0cm}}
\toprule
Variable & BL & FU & Values & Role & Note \\
\midrule
\endhead
\var{ct_cag_<v>_sis} & \yes & \yes & Count of segments with plaque & check &
  Agrees with the segment reads in 95--97\,\% for CX and RCA but only 64--79\,\% for LAD; the recorded LAD
  score may include LM or IM. \\
\var{ct_cag_<v>_sss} & \yes & \yes & Sum of stenosis grades & check & As above. \\
\var{ct_cag_<v>_plaque_type_score} & \yes & \yes & Sum of plaque-type scores & not used & \\
\var{ct_cag_<v>_uneval} & \yes & \yes & Count of unevaluable segments & check & \\
\var{ct_cag_<v>_uneval_contrast}, \var{_scanfield}, \var{_pm}, \var{_noise} & \no & \yes & Count & check &
  Reasons for unevaluable segments at FU. \\
\var{ct_cacs_<v>} (\var{lm}, \var{lad}, \var{cx}, \var{rca}) & \yes & \yes & Agatston score & check, covariate &
  Calcium per vessel; secondary outcome (0 $\rightarrow$ $>$0) and combined eligibility (trunk none and vessel calcium 0). \\
\bottomrule
\end{longtable}

\section{Patient-level summaries}

\begin{longtable}{L{4.3cm} c c L{6.2cm} L{2.2cm} L{7.0cm}}
\toprule
Variable & BL & FU & Values & Role & Note \\
\midrule
\endhead
\var{ct_cag_sis} & \yes & \yes & Segment involvement score & covariate &
  Recorded as 0, not missing, for the 1\,094 patients without BL segment reads. \\
\var{ct_cag_sss} & \yes & \yes & Segment stenosis score & covariate & As above. \\
\var{ct_cag_plaque_type_score} & \yes & \yes & Sum of plaque-type scores & not used & \\
\var{ct_cag_none} \dots\ \var{ct_cag_occlusion} & \yes & \yes & Flags or counts per stenosis category & not used &
  Six variables (none, minimal, mild, moderate, severe, occlusion); all 0 for patients without BL reads. \\
\var{ct_cag_dominance} & \yes & \yes & Text & covariate & Needed to assign PDA and PLA. \\
\var{ct_cag_stent}, \var{ct_cag_cabg} & \yes & \yes & 0/1 & eligibility & Prior revascularisation; stent missing for $\approx$66\,\% at BL. \\
\var{ct_cag_incomplete} & \yes & \yes & Count & check & 0 for the 1\,022 read patients and 16--17 for the others at BL. \\
\var{ct_cag_uneval} & \yes & \yes & Count of unevaluable segments & check & \\
\var{ct_cag_uneval_motion}, \var{_noise}, \var{_scanfield}, \var{_contrastdefect}, \var{_pm}, \var{_other} & \no & \yes & Count & check &
  Reasons at FU; motion dominates. \\
\var{ct_cag_eval_n} & \yes & \no & Number of evaluable segments & check & Median 16. \\
\var{ct_cacs_sum} & \yes & \yes & Total Agatston score & covariate, check &
  Available for 2\,109 patients at both visits; calcium progression (Hokanson) is the most reliable outcome in the data. \\
\var{ct_cacs_groups} & \yes & \yes & Category 1--3 & not used & \\
\var{ct_cacs_<v>_eval}, \var{ct_cacs_eval_all} & \no & \yes & 0/1 & check & Calcium evaluability at FU. \\
\bottomrule
\end{longtable}

\section{Clinical events and timing}

\begin{longtable}{L{4.3cm} c c L{6.2cm} L{2.2cm} L{7.0cm}}
\toprule
Variable & BL & FU & Values & Role & Note \\
\midrule
\endhead
\var{scandate} & \yes & \yes & Date & covariate &
  Gives the follow-up time (median 10.0 years, 5th--95th percentile 9.5--15.3), used as log offset. \\
\var{demo_hist_myocardial_infarction} & \yes & \yes & 0/1 & not used & 6 new events; missing for $\approx$51\,\% at BL. \\
\var{demo_hist_coronary_revasc} & \yes & \yes & 0/1 & not used & 6 new events. \\
\var{demo_hist_stroke_tia} & \yes & \yes & 0/1 & not used & 31 new; 119 report it at BL but not at FU. \\
\var{hist_angina} & \yes & \yes & 0/1 & not used & 96 new. \\
\var{demo_hist_afib}, \var{_heart_failure}, \var{_pad}, \var{_icd_pm}, \var{_cag} & \no & \yes & 0/1 & not used &
  FU only; new cannot be separated from pre-existing. \\
\var{ae_*} & \no & \yes & Trial adverse events & not used & Related to the trial, not to coronary disease. \\
\bottomrule
\end{longtable}

\noindent Hard clinical events (myocardial infarction, revascularisation, stent, bypass) total 10 over follow-up
and cannot serve as an outcome.

\section{Image quality}

\begin{longtable}{L{4.3cm} c c L{6.2cm} L{2.2cm} L{7.0cm}}
\toprule
Variable & BL & FU & Values & Role & Note \\
\midrule
\endhead
\var{scan_qualitative_quality} & \yes & \no & 0/1 & check & Missing for 246. \\
\var{bad_quality} & \yes & \no & True/False & not used & False for every patient. \\
\var{Image_Quality_reason_1}, \var{_2} & \yes & \no & Free text & not used & Not reported by the script (text). \\
\var{ct_angio_heart_done}, \var{ct_noncontrast_heart_done} & \no & \yes & 0/1 & check &
  1 for every patient at FU; no BL counterpart, so BL contrast CT availability must be checked from the volumes. \\
\bottomrule
\end{longtable}

\section{Chosen design and derived outcomes}
\label{sec:design}

\begin{longtable}{L{4.3cm} L{19cm}}
\toprule
Element & Definition \\
\midrule
\endhead
Unit & LAD (\var{plad}, \var{mlad}, \var{dlad}), CX (\var{pcx}, \var{mcx}, \var{dcx}) and RCA (\var{prca}, \var{mrca}, \var{drca}) trunks; side branches excluded. \\
Eligibility & All trunk segments evaluable and graded none at BL; for LAD and CX the LM is also graded none. \\
Primary outcome & Any trunk segment graded minimal or worse at FU (binary). Unknown if no plaque is seen and a trunk segment is unevaluable at FU. \\
Secondary outcome & Worst FU grade in the trunk, ordinal: none, minimal, mild or worse. \\
Calcium check & Vessel calcium 0 at BL $\rightarrow$ $>$0 at FU (\var{ct_cacs_lad}, \var{ct_cacs_cx}, \var{ct_cacs_rca}). \\
Sensitivity & Eligibility additionally requiring vessel calcium 0 at BL; RCA additionally requiring LM none. \\
Follow-up time & $\log$(years between \var{scandate} values) as offset with a complementary log-log link. \\
\bottomrule
\end{longtable}

\subsection*{Open issues}
\begin{itemize}
  \item Whether the 1\,094 patients without BL segment reads had a normal BL scan or no contrast CT;
        decides whether they can enter the design.
  \item The CGPS stenosis grading thresholds and segment boundary definitions (reading protocol).
  \item The meaning of the \var{ct_cag_<s>_eval} codes and of \var{ct_cag_incomplete}.
\end{itemize}

\begin{thebibliography}{9}
\bibitem{leipsic2014} J.~Leipsic et al. SCCT guidelines for the interpretation and reporting of coronary CT angiography.
  \emph{Journal of Cardiovascular Computed Tomography}, 8(5):342--358, 2014.
\bibitem{cury2016} R.~C.~Cury et al. CAD-RADS: Coronary Artery Disease -- Reporting and Data System.
  \emph{Journal of Cardiovascular Computed Tomography}, 10(4):269--281, 2016.
\end{thebibliography}

\end{document}
